\chapter{Build and Understand the Bot}

\begin{goalbox}
Connect the token to TeleBot and understand how each message gets its reply.
\end{goalbox}

The repository already contains the completed file below. Reading it from top to
bottom is the fastest way to connect each idea to working code.

\lstinputlisting[language=Python,caption={The complete code in code/bot.py}]{\CodePath/bot.py}

\section{Step 1: import the tools}

\texttt{os} reads the environment variable. \texttt{telebot} creates the client
that talks to Telegram. \texttt{types} supplies message type hints; these hints
help readers and editors but do not change the bot's behavior.

\section{Step 2: keep replies easy to test}

The functions \texttt{start\_reply}, \texttt{help\_reply}, and
\texttt{echo\_reply} only build text. They do not need the network. Separating
them makes the expected conversation clear and lets the offline tests check it.

\section{Step 3: register handlers}

A \emph{handler} is simply a function chosen for a kind of incoming message.
The decorators above the functions define the choice:

\begin{itemize}
\item \texttt{commands=["start"]} handles \texttt{/start};
\item \texttt{commands=["help"]} handles \texttt{/help};
\item \texttt{content\_types=["text"]} handles other text;
\item the final handler covers common non-text messages.
\end{itemize}

Order matters: command handlers appear before the general text handler, so a
command receives its special answer.

\section{Step 4: start polling}

\texttt{main} stops with a readable error when \texttt{BOT\_TOKEN} is absent.
When it exists, \texttt{infinity\_}\allowbreak\texttt{polling} keeps listening until you press
\texttt{Ctrl+C}. The \texttt{skip\_}\allowbreak\texttt{pending=True} option ignores old messages
that accumulated while this learning bot was offline.

\begin{checkpointbox}
You can point to the exact handler used for \texttt{/start}, \texttt{/help}, a
normal text message, and a photo.
\end{checkpointbox}
